% Source: https://artofproblemsolving.com/community/c3h1748217_geo_question
% Stable identity: aops-topic-1748217-q01
Let $r$ be the radius and let the chords of lengths $10$ and $8$ be at distances $d$ and $e$ from the centre $O$. The perpendicular from $O$ to a chord bisects it, so
\[
r^2=5^2+d^2=4^2+e^2,\qquad\text{hence}\qquad e^2-d^2=9 .
\]
\textbf{Same side of the centre.} Then $e-d=6$, so $e+d=\dfrac{9}{6}=1.5<6$, which would make $d$ negative. Impossible.

\textbf{Opposite sides.} Then $d+e=6$ and $e-d=\dfrac{9}{6}=1.5$, so $d=2.25$ and $e=3.75$.

\begin{asy}
usepackage("amsmath");
size(5.5cm);
defaultpen(fontsize(10pt));
real r=sqrt(481)/4;pair O=(0,0);
pair A1=(-5,2.25),A2=(5,2.25),B1=(-4,-3.75),B2=(4,-3.75);
draw(circle(O,r),linewidth(0.8));
draw(A1--A2,linewidth(0.9));draw(B1--B2,linewidth(0.9));
draw((0,2.25)--(0,-3.75),dashed);
draw(O--A2,blue);draw(O--B2,blue);dot(O);
draw((0,2.25)--(0.3,2.25)--(0.3,1.95)--(0,1.95),linewidth(0.5));
draw((0,-3.75)--(0.3,-3.75)--(0.3,-3.45)--(0,-3.45),linewidth(0.5));
label("$O$",O,W);
label("$5$",(2.5,2.25),N);label("$4$",(2,-3.75),S);
label("$d$",(0,1.1),W);label("$6-d$",(0,-1.9),W);
label("$r$",(A2+O)/2+(0,0.1),N,blue);label("$r$",(B2+O)/2+(0.25,0),E,blue);
label("\mbox{}",(min(currentpicture,user=true)+max(currentpicture,user=true))/2,black); // keeps the text after the figure black when LaTeX typesets the labels inline
\end{asy}

Hence $r^2=25+2.25^2=25+\dfrac{81}{16}=\dfrac{481}{16}$ (check: $16+3.75^2=16+\dfrac{225}{16}=\dfrac{481}{16}$), and
\[
r=\boxed{\dfrac{\sqrt{481}}{4}}\approx5.48 .
\]
